% \begin{savequote}
% \begin{verbatim}
%  ________________________________________
% / You mentioned your name as if I should \
% | recognize it, but beyond the obvious   |
% | facts that you are a bachelor, a       |
% | solicitor, a freemason, and an         |
% | asthmatic, I know nothing whatever     |
% | about you.                             |
% |                                        |
% | -- Sherlock Holmes, "The Norwood       |
% \ Builder"                               /
%  ----------------------------------------
%        \   ,__,
%         \  (oo)____
%            (__)    )\
%               ||--|| *
% 
% \end{verbatim}
% \end{savequote}
\def\ChapterCode{TOD}
\chapter
[\CO{[TOD]} Sterben und Tod]
{Sterben und Tod}
\label{chp:TOD}

\ChapterIntro
{%
\noindent Das Sterben und der Tod sind alltägliche Begleiter
des Lebens. Im Gesundheitsbereich ist man zwangsläufig
regelmäßig auf die eine oder andere Weise damit
konfrontiert. Ziel ist es, den legitimen Wünschen der
Patienten und Angehörigen Rechnung zu tragen, und
andererseits selbst mit zum Teil sehr fordernden Situationen
und Fragen zurecht zu kommen.
}
{%
\begin{description}
  \item[Maintainer:] Sebastian Gabriel
  \item[Autoren:] Diverse
  \item[Reviewer:] Standard-Reviewprozess
  \item[Version:] {\VarDocVersionTypeName} ({\VarDocVersionTypeDescription})
  \item[SHA1:] {\DocGitCommit}
\end{description}

{\TxTeilkapitelAASS}
}

\Definition{Der \T{Tod} ist das Ende der Lebensfunktionen
und der Zustand danach. Das \T{Sterben} ist ein mehr oder
weniger lang dauernder Prozess, an dessen Ende der Tod
steht.}

\section{Der Tod allgemein}
\index{Tod}
\label{topic:tod}


Der Tod gehört zu den einschneidensten und
unausweichlichsten Dingen des Lebens. 
Man
unterscheidet zwischen dem klinischen und dem biologischen
Tod

% \paragraph{Klinischer Tod vs. biologischer Tod}
\index{Tod!klinischer}\index{Tod!biologischer}

\begin{center}
\CorrectionItemizeBox\FontTableBase
\begin{tabularx}{\linewidth}{XX}
\toprule\addlinespace
\noindent\TabHeading{Klinischer Tod}
 &
\noindent\TabHeading{Biologischer Tod}
\\\addlinespace\midrule\addlinespace
Keine Atmung

Kein Puls

Unter Umständen reversibel (Reanimation)
 &
Organtod des Gehirns (irreversibel)

Gleichzusetzen mit Tod des Individuums
\\\addlinespace\bottomrule
\end{tabularx}
\end{center}

\section{Todesfeststellung}


\AuthNotes{Todeszeichen: Nicht wie von Roman gefordert in
Tabelle mit mit sicheren und unsicheren Todeszeichen
umgebaut, halte ich für unübersichtlich und befürchte
Verwechslungsgefahr mit Knochenbruchzeichen.[GAB]}





\Text{Zeichen des Todes und Leichenzeichen}
{%
Grundsätzlich gibt es Zeichen, die als \EE{sichere
Todeszeichen}
gelten~\cite{ForMed2,GrassbergerTodesermittlung1}:

\begin{enumerate}
    \item \EE{Bildung von blauvioletten Totenflecken} durch
    Absackung des Blutes an die untenliegenden Stellen.
    \item \EE{Totenstarre}
    \item \EE{Nicht mit dem Leben vereinbare Verletzungen}
    wie zum Beispiel eine Abtrennung des Kopfes vom Körper.
    \item \EE{Fäulniserscheinungen}
\end{enumerate}

\noindent Die ersten drei Punkte sind \E{frühe
Leichenerscheinungen}, da sie relativ rasch auftreten.
Andere frühe Leichenerscheinungen sind die \E{Abkühlung der
Leiche} und die \E{Totenstarre}. Längere Zeit nach dem Tod
treten die \E{späten Leichenerscheinungen} auf, dazu gehören
\E{Fäulnis und Verwesung}, \E{Mumifizierung}, \E{Tierfraß}
und \E{Skelettierung}.

Die Feststellung der Todeszeichen und Leichenerscheinungen
ist jedoch oft nicht so einfach, die Anzeichen können oft
auch mit anderen Sachen verwechselt werden, z.\,B.
Totenflecke mit Blutergüssen oder die Totenstarre mit einem
tonischen Krampf.
}{}{}{}{}

\Fazit{Es gibt sichere Todeszeichen, das Erkennen dieser
Zeichen ist aber unsicher und bedarf Wissen und Erfahrung!}

\Fazit{Die Todesfeststellung erfolgt grundsätzlich nur durch einen
Arzt!}\index{Todesfeststellung}

\Cave{Für unterkühlte Patienten gilt: 
\Speech{\enquote{Nobody is dead until warm and dead!}}

(Niemand wird für tot erklärt, bevor er nicht vorher erwärmt
wurde.)}

%\vspace{1em}
%%%%%%%%%%%%%%%%%%%%%%%%%%%%%%%%%%%%%%%%%%%%%%%%%%%%%%%%%%%
%%%%%%%%%%%%%%%%%%%%%%%%%%%%%%%%%%%%%%%%%%%%%%%%%%%%%%%%%%%
%%%%%%%%%%%%%%%%%%%%%%%%%%%%%%%%%%%%%%%%%%%%%%%%%%%%%%%%%%%
\section{Unterlassung einer Reanimation}
\label{topic:reanimation-unterlassung}



\MassnahmenMK{Spezielle Maßnahmen}{Unterlassung der Reanimation}
{m:reanimation-unterlassung}
{%
Unter bestimmten Umständen darf
(oder muss) eine Reanimation unterlassen
werden.\index{Reanimation!Unterlassung
der}\index{Wiederbelebung|see{Reanimation}}

}{}
% \MassnahmenWide{Spezielle Maßnahmen}{Unterlassung der Reanimation}
% {\label{m:reanimation-unterlassung}}
% {%
% \fcolorbox{ubuntured}{White}{
% \begin{minipage}[t]{\linewidth - {4\fboxrule} - {4\fboxsep}}
% \Ca{Unter den folgenden Voraussetzungen darf (oder muss) eine
% Reanimation unterlassen
% werden:\index{Reanimation!Unterlassung
% der}\index{Wiederbelebung|see{Reanimation}}}
% \begin{multicols}{2}              
% \begin{itemize}
%   \item \EE{Verwesung}
%   \item \EE{Fäulnis}
%   \item \EE{Mumifizierung}
%   \item \EE{Skelettierung}
%   \item \EE{Verletzungen, die nicht mit dem Leben
%   vereinbar sind}
%   \item \EE{Tierfraß}\footnote{Unter
%   \E{\enquote{Tierfraß}} versteht man den Verzehr eines Kadavers
%   durch Kadaverfresser. Ein frischer Hundebiss oder
%   Ähnliches fällt \underbar{nicht} darunter. Maden können
%   sich im Gewebe von Lebenden einnisten und sind daher
%   auch kein sicheres Todeszeichen.}
%   \item (Verbindliche) \EE{Patientenverfügung}\footnote{\E{Patientenverfügung}: Bei Notfallversorgungen darf \EE{nicht nach einer Verfügung gesucht
%   werden, wenn dadurch das Leben oder die Gesundheit des Patienten ernstlich
%   gefährdet werden würde}.
%   \EE{Aber:} Wenn \E{zweifelsfrei} eine verbindliche
%   Patientenverfügung vorliegt (z.\,B. weil eine betreuende
%   Pflegekraft diese in der Zwischenzeit gefunden und die
%   Verbindlichkeit bestätigt hat), muss dieser Folge
%   geleistet werden.
% Liegt eine \E{beachtliche Patientenverfügung} vor, so ist im
% \E{Einzelfall} zu entscheiden: Das Wohl des Patienten bleibt
% oberstes Gebot und die Verfügung soll in die
% Entscheidung einfließen.}\index{Patientenverfügung!Unterlassung der
%   Reanimation} (\FullPrettyRef{topic:patientenverfuegung}; 
%   Bis zur zweifelsfreien Klärung der Situation muss
%   jedenfalls eine Reanimation begonnen bzw. fortgeführt
%   werden!)
%   \item (Abbruch der Reanimationsmaßnahmen aufgrund
%   von \EE{Erschöpfung})\footnote{Erschöpfung darf im
%   professionellen Umfeld keine Rolle spielen, für eine
%   entsprechende Ablösung ist bei Bedarf zu sorgen.}
%   \item Unvereinbarkeit mit
%   \EE{Selbstschutz}\footnote{Es müssen jedoch alle Maßnahmen
%   getroffen werden, die ohne Gefährdung des Selbstschutzes
%   möglich sind. So wird zum Beispiel für Ersthelfer
%   empfohlen, wenn
%   eine Mund-zu-Mund-Beatmung nicht zumutbar ist, zumindest 
%   eine Herzdruckmassage durchzuführen.}
% \end{itemize}
% 
% \end{multicols}
% \end{minipage}
% }
% }



\section{Begleitung und Betreuung Sterbender}

\label{topic:sterbende-betreuung}



Der Tod ist ein integraler Bestandteil des Lebens. So banal
diese Aussage auch scheinen mag, so ambivalent ist das
Verhalten des bzw. der Menschen zu ihm. Nach tausenden
Jahren stellt der Tod und der Sterbeprozess noch immer
(oder vielleicht sogar mehr denn je) eine große
Herausforderung sowohl für den Sterbenden, als auch für
dessen Umgebung dar. 

% \expandedtitle{simple}{\Speech{\enquote{Der Tod ist etwas
% persönliches.}}}




Die Auseinandersetzung mit dem Tod und der eigenen
Endlichkeit verläuft bei jedem Menschen äußerst
unterschiedlich und ist von dessen Weltanschauung, Umgebung,
Lebensgeschichte und -situation und vielen anderen Faktoren
abhängig, und kann im Verlauf des Lebens ebenso radikalen
Änderungen unterworfen sein. Mögliche Bewältigungstrategien
sind:

\begin{itemize}
  \item Verdrängung
  \item Hoffnungen, welche sich an den Tod anknüpfen
  (Wiedergeburt, Leben im Jenseits)
  \item Vermächtnis an die Nachwelt; Schaffen von etwas, was
 \enquote{größer als man selbst ist}
  \item Rationalisierung (Vorsorgeversicherungen,
  Testamente, \ldots)
  \item Rituale (Gebete, Begräbniszeremonien,
  Leichenschmaus, \ldots)
\end{itemize}

\Fazit{Für den Umgang mit Sterben und dem Tod gibt es kein
Patentrezept und keinen\enquote{`goldenen Weg}.}








Elisabeth \T{Kübler-Ross} beschreibt fünf Phasen, die im
Sterbeprozess in unterschiedlicher Reihenfolge, oft auch
wiederholt, durchlaufen werden:

\begin{itemize}
  \item \EE{Abwehr}: Nicht wahrhaben wollen, Abwehr,
  Verdrängung
  \item \EE{Zorn}: Wütend auf das
  Schicksal (\Speech{\enquote{Warum ausgerechnet ich und nicht
  ein anderer?}}) und neidig auf das \enquote{Glück} anderer
  \item \EE{Verhandeln} {\ldots} mit dem
  betreuenden Personal, mit Gott oder generell mit dem
  Schicksal
  \item \EE{Depression}: Traurige,
  deprimierte Stimmungslage, Verlust des \enquote{sinnerfüllten
  Daseins}; vgl. \FullPrettyRef{topic:depression}
  \item \EE{Akzeptanz}: Der Patient kann sein Schicksal
  annehmen, er findet sich damit ab.
\end{itemize}

Demnach ist es verständlich, dass der gleiche Patient zu
unterschiedlichen Zeitpunkten völlig unterschiedlich
reagieren kann. Was auf den ersten Blick wie \enquote{launisches
Verhalten} wirkt, ist oft durch die Phasen des
Sterbeprozesses erklärbar.


\Text{Palliativbehandlung}
{Die \T{Palliativbehandlung} bzw. -betreuung beschäftigt
sich mit der Betreuung von Patienten, für die es keine heilende Therapie mehr
gibt, oder sich gegen eine solche entschieden haben.
Ziel ist nicht die Heilung einer Krankheit, sondern das
Ermöglichen eines würdigen und schmerzfreien Lebens bis zum Eintreten des Todes.
Die Bedürfnisse der Patienten sind unterschiedlich und richten
sich zum Beispiel nach dem Alter, der Erkrankung, den
familiären Umständen, der Religion, usw.}
{}
{}
{}
{}


\Text{Der Umgang mit \enquote{geplant} sterbenden Patienten im
Rettungs- und Krankentransportdienst} {Die Zeitdauer des
Kontakts ist naturgemäß eher gering, und die palliative
Versorgung ist nicht unbedingt ein Schwerpunkt während eines
Transportes.  Dennoch gibt es einige Dinge zu bedenken:

\begin{itemize}
  \item \EE{Angst und Ausnahmesituation}: Sowohl die Angehörigen wie
  auch der Patient befinden sich u.\,U. in einer
  Ausnahmesituation wenn eine (neuerliche) Klinikeinweisung
  bevorsteht. Oft kann man Angst vor der folgenden
  Ungewissheit beobachten. Hier ist es wichtig,
  selber ein \E{ruhiger Pol} bzw. ein \E{Fels in der Brandung}
  zu sein, an dem sich der Patient und die Angehörigen
  \enquote{anhalten} können.
  \item \EE{Sachebene als Brücke}: Sachliche und klare Informationen
  über das weitere Vorgehen sind für Patienten und
  Angehörige wichtig. Sie geben ihnen (und auch dem
  Personal) Sicherheit im Gespräch und können ein erster
  Schritt hin zu einer vertrauensvollen
  Gesprächsbasis sein. Es kann damit auch gelingen, so
  manche überschießende emotionale Reaktion zu kanalisieren.
  \item \EE{Verweigerung der Einweisung}: Viele Patienten wollen
  nicht in einem Spital oder einer anderen Institution,
  sondern zu Hause sterben. Dieser Wunsch ist zu akzeptieren.
  Es muss abgeklärt werden, ob der Patient einsichts-
  und urteilsfähig (Bedenke: Besachwaltung!) und wie er sonst
  versorgt (Essen auf Rädern, Heimhilfe, Angehörige, \ldots) ist.
  \item \EE{Patientenverfügungen}: Eventuell vorhandene
  Patientenverfügungen (verbindliche und beachtliche;
  \FullPrettyRef{topic:patientenverfuegung}) sollen, wenn
  vorhanden, mitgenommen werden.
  \item \EE{Eigene Psychohygiene}
\end{itemize}
}
{}
{}
{}
{}






\Text{Der Helfer und der Sterbende}
{Der Kontakt und der Umgang mit Sterbenden ist oft auch für
den Helfer problematisch und kann eine psychische Belastung
darstellen. Das Besondere am Umgang mit
Sterbenden ist, dass viele Dinge, die sonst im Leben und in
der Ausübung von Gesundheitsberufen selbstverständlich sind,
hier keine Geltung haben. Der Tod soll dort normalerweise um
jeden Preis vermieden und bekämpft werden. Derartige, in
diesem Bereich unerfüllbare, Zielsetzungen (\enquote{Heilen
wollen}) führen zwangsläufig zu einer Überforderung.

Folgende Leitsätze können unter Umständen hilfreich sein:

\begin{itemize}
  \item \EE{Persönliche Abgrenzung}: Das Problem des
  Patienten bleibt grundsätzlich das Problem des Patienten
  \item Der Helfer muss im Rahmen der Möglichkeiten im Sinne
  des Patienten handeln und seine angemessenen
  Bedürfnisse befriedigen.
%   \item Wenn der Helfer mit unzufrieeden mit der Betreuung
%   ist, muss er versuchen, den Rahmen des Möglichen
%   zu erweitern.
\item Man darf nicht seine eigenen Vorstellungen
  bezüglich der Hilfe mit denen des Patienten
  gleichstellen. Wenn man dem Patienten helfen möchte, muss
  man wissen, was \E{er} tatsächlich will.
\end{itemize}

\Fazit{Man
muss sich also ständig vor Augen halten, dass im Umgang mit
Sterbenden manche, sonst übliche, Ansichten nicht -- oder
nur eingschränkt -- gelten.}
}
{}
{}
{}
{}





% \expandedtitle{roundbox}{\Speech{\enquote{Das große Problem hat das
% Fachpersonal.}}}





% %%%%%%%%%%%%%%%%%%%%%%%%%%%%%%%%%%%%%%%%%%%%%%%%%%%%%%%%%%%
% %%%%%%%%%%%%%%%%%%%%%%%%%%%%%%%%%%%%%%%%%%%%%%%%%%%%%%%%%%%
% %%%%%%%%%%%%%%%%%%%%%%%%%%%%%%%%%%%%%%%%%%%%%%%%%%%%%%%%%%%
% %%%%%%%%%%%%%%%%%%%%%%%%%%%%%%%%%%%%%%%%%%%%%%%%%%%%%%%%%%%
% \subsection{Der unerwartete Tod}













\AuthNotes{GABG: gesicherte bösartige Erkrankung?

GABS: Was gilt als gesichert? Was gilt als bösartig? Eine
bösartige Grunderkrankung gilt nicht per se als Argument um
nicht zu reanimieren $\rightarrow$ Patientenverfügung,
mutmaßlicher Patientenwille. Prostatakarzinom und Infarkt?}



% \section*{\protect\dsliterary{} Weiterführende Literatur}
% \begin{WidePar}
% \begin{multicols}{2}
% \Z{%
% \begin{labeling}{mmm}
%   \item[\citep{NagelePalliativ1}] \emph{\bibentry{NagelePalliativ1}} 
%   \item[\citep{ForMed2}] \emph{\bibentry{ForMed2}}
% %   \item[\citep{ScholzNotfallmedizin2}] \emph{\bibentry{ScholzNotfallmedizin2}} 
% %   \item[\citep{IntroClinEmergMed4}] \emph{\bibentry{IntroClinEmergMed4}} 
% %   \item[\citep{Longmore2007}] \emph{\bibentry{Longmore2007}} 
% %   \item[\citep{Fauci2008}] \emph{\bibentry{Fauci2008}}
% \end{labeling}
% }
% \end{multicols}
% \end{WidePar}

\citep{NagelePalliativ1,ForMed2}

%O2:
%\nocite{Morris2001,Milross1989,Boumphrey2003}
%\nocite{BoehmerTaschRett6,Fauci2008,LPN3,LippertAnatomie7,StriebelAn7,DRAn3}


%\putbib

\ChapterBibliography



\ChapterEnd
